\subsection{PRELIMINARIES TO THE CONSTRUCTION OF ROOT SYSTEMS}

Let V be a real vector space of dimension \(l \geqslant 1\) equipped with a scalar product \((x|y)\) . L a discrete subgroup of V, A a finite set of numbers \textgreater0, and R the set of \(\alpha \in L\) such that \((\alpha|\alpha) \in \Lambda\) . Assume that R generates V and that, for any pair \((\alpha,\beta)\) of points of R, the number \(2\frac{(\alpha|\beta)}{(\alpha|\alpha)}\) is an integer. Then, R is a root system in V. Indeed, R clearly satisfies \((RS_{I})\) . Let \(\alpha \in R\) : let \(s_{\alpha}\) be the orthogonal reflection \(x \mapsto x - 2\frac{(x|\alpha)}{(\alpha|\alpha)}\alpha\) ; then, if \(\beta \in R\) , we have \(2\frac{(\beta|\alpha)}{(\alpha|\alpha)} \in Z\) , so \(s_{\alpha}(\beta) \in L\) , and moreover \(\|s_{\alpha}(\beta)\| = \| \beta \|\) , so \(s_{\alpha}(\beta) \in R\) ; thus, R satisfies \((RS_{II})\) and \((RS_{III})\) , and is reduced if A does not contain two numbers of the form \(\lambda\) and \(4\lambda\) .

Now let V be a subspace of \(E = R^{n}\) . Let \((\varepsilon_{1}, \ldots, \varepsilon_{n})\) be the canonical basis of E; we equip E with the scalar product for which this basis is orthonormal and identify \(E^{*}\) with E (resp. \(V^{*}\) with V) by means of this scalar product. We define subgroups \(L_{0}, L_{1}, L_{2}, L_{3}\) of E as follows:

\begin{enumerate}
\def\labelenumi{\arabic{enumi})}
\item
  \(L_{0}\) is the Z-module with basis \((\varepsilon_{i})\) . We have \((\alpha|\beta)\in\mathbf{Z}\) for all \(\alpha,\beta\in L_{0}\) . The vectors \(\alpha\in L_{0}\) for which \((\alpha|\alpha)=1\) are the \(\pm\varepsilon_{i}\;(1\leqslant i\leqslant n)\) ; those for which \((\alpha|\alpha)=2\) are the \(\pm\varepsilon_{i}\pm\varepsilon_{j}\) for i\textless j (the two signs \(\pm\) , in \(\pm\varepsilon_{i}\pm\varepsilon_{j}\) , are chosen independently of each other; we adopt an analogous convention throughout the remainder of this section).
\item
  \(L_{1}\) is the Z-submodule of \(L_{0}\) consisting of the \(x = \sum_{i=1}^{n} \xi_{i} \varepsilon_{i} \in L_{0}\) such that \(\sum_{i=1}^{n} \xi_{i}\) is even; since \(\xi_{i}\) and \(\xi_{i}^{2}\) have the same parity, this is equivalent to requiring that \((x|x)\) is even. Let \(L_{1}^{\prime}\) be the submodule of \(L_{1}\) generated by the \(\varepsilon_{i} \pm \varepsilon_{j}\) ; we have \(\sum_{i=1}^{n} \xi_{i} \varepsilon_{i} \equiv (\sum_{i=1}^{n} \xi_{i}) \varepsilon_{n} \pmod{L_{1}^{\prime}}\) , and since \(2\varepsilon_{n} = (\varepsilon_{1} + \varepsilon_{n}) - (\varepsilon_{1} - \varepsilon_{n}) \in L_{1}^{\prime}\) , it follows that \(L_{1}^{\prime} = L_{1}\) . Since \(L_{0}\) is generated by \(L_{1}\) and \(\varepsilon_{1}, L_{0}/L_{1}\) is isomorphic to Z/2Z.
\item
  \(\mathrm{L}_2 = \mathrm{L}_0 + \mathbf{Z}\left(\frac{1}{2} \sum_{i=1}^{n} \varepsilon_i\right)\) . It is clear that an element \(x = \sum_{i=1}^{n} \xi_i \varepsilon_i\) of \(\mathrm{V}\) is in \(\mathrm{L}_2\) if and only if
\end{enumerate}

\[
2 \xi_ {i} \in \mathbf {Z}, \quad \xi_ {i} - \xi_ {j} \in \mathbf {Z} \text {   for   all   } i \text {   and   } j.\tag{6}
\]

Since \((\varepsilon_k|_{\frac{1}{2}}\sum_{i=1}^{n}\varepsilon_i) = \frac{1}{2}\) for all \(k\) , and since \(\| \frac{1}{2}\sum_{i=1}^{n}\varepsilon_i\|^2 = \frac{n}{4}\) , we have \((\alpha|\beta) \in \frac{1}{2}\mathbf{Z}\) for \(\alpha, \beta \in L_2\) if \(n\) is even. The group \(L_2 / L_0\) is isomorphic to \(\mathbf{Z}/2\mathbf{Z}\) .

\begin{enumerate}
\def\labelenumi{\arabic{enumi})}
\setcounter{enumi}{3}
\tightlist
\item
  \(L_{3} = L_{1} + \mathbf{Z}\left(\frac{1}{2} \sum_{i=1}^{n} \varepsilon_{i}\right)\) . If n is a multiple of 4, \(L_{3}\) is the set of \(\sum_{i=1}^{n} \xi_{i} \varepsilon_{i}\) satisfying (6) and the condition \(\sum_{i=1}^{n} \xi_{i} \in 2\mathbf{Z}\) : in this case, \((\alpha | \beta) \in \mathbf{Z}\) for all \(\alpha, \beta \in L_{3}\) .
\end{enumerate}

It is immediate that the subgroup of E associated to \(L_{0}\) (resp. \(L_{1}, L_{2}\) ) is \(L_{0}\) (resp. \(L_{2}, L_{1}\) ). The subgroup of E associated to \(L_{3}\) is the set of

\[
x = \sum_ {i = 1} ^ {n} \xi_ {i} \varepsilon_ {i} \in L _ {2}
\]

such that \((x|\frac{1}{2}\sum_{i = 1}^{n}\varepsilon_i)\in \mathbf{Z}\) , that is, such that \(\sum_{i = 1}^{n}\xi_{i}\in 2\mathbf{Z}\) : if \(n\equiv 0\) (mod. 4), this associated subgroup is therefore \(\mathrm{L}_3\) .

The abelian group \(\mathrm{L}_2 / \mathrm{L}_1\) is of order 4, and hence is isomorphic to \(\mathbf{Z} / 4\mathbf{Z}\) or \(\mathbf{Z} / 2\mathbf{Z} \times \mathbf{Z} / 2\mathbf{Z}\) (Algebra, Chap. VII, § 4, no. 6, Th. 3). If \(n\) is odd,

\[
p (\frac {1}{2} \sum_ {i = 1} ^ {n} \varepsilon_ {i}) \in \mathrm{L} _ {1} \Leftrightarrow p \equiv 0 (\mathrm{mod.} 4)
\]

so \(\mathrm{L}_2 / \mathrm{L}_1\) is cyclic of order 4. If \(n\) is even,

\[
p (\frac {1}{2} \sum_ {i = 1} ^ {n} \varepsilon_ {i}) \in \mathrm{L} _ {1} \Leftrightarrow p \equiv 0 (\mathrm{mod.} 2)
\]

so \(L_{2}/L_{1}\) , which contains two distinct elements of order 2, is isomorphic to \(Z/2Z \times Z/2Z\) .

We shall use this notation throughout the next nine nos. and in the plates. For each type of Dynkin graph in Th. 3, we shall describe explicitly:

\begin{enumerate}
\def\labelenumi{(\Roman{enumi})}
\item
  A root system R and the number of roots.
\item
  A basis B of R. and the corresponding positive roots. The basis B will be indexed by the integers \(1,\ldots,l\) .
\item
  The Coxeter number \(h\) (§ 1. no. 11).
\item
  The highest root \(\tilde{\alpha}\) (for the order defined by B) and the completed Dynkin graph (no. 3). We shall indicate next to each vertex the corresponding root of B.
\item
  The inverse system \(\mathbb{R}^{\prime}\) , the canonical bilinear form and the constant \(\gamma(\mathbb{R})\) (§ 1, no. 12).
\item
  The fundamental weights relative to B (§ 1, no. 10).
\item
  The sum of the positive roots.
\item
  The groups \(P(R), Q(R), P(R)/Q(R)\) and the connection index (§ 1, no. 9).
\item
  The exponents of \(\mathrm{W}(\mathrm{R})\) (Chap. V. § 6, no. 2, Def. 2). In the cases \(\mathbf{A}_l, \mathbf{B}_l, \mathbf{C}_l\) and \(\mathbf{D}_l\) we determine the symmetric invariants.
\item
  The order of W(R) (and in some cases its structure).
\item
  The group \(A(R)/W(R)\) , its action on the Dynkin graph, and the element \(w_{0}\) of \(W(R)\) that transforms B to -B.
\item
  The action of \(\mathrm{P}(\mathrm{R}^{\prime})/\mathrm{Q}(\mathrm{R}^{\prime})\) on the completed Dynkin graph and the action of \(\mathrm{A}(\mathrm{R})/\mathrm{W}(\mathrm{R})\) on \(\mathrm{P}(\mathrm{R}^{\prime})/\mathrm{Q}(\mathrm{R}^{\prime})\) .
\end{enumerate}

For each Dynkin graph in Th. 3, these data will be collected in Plates I to IX, and ordered in a uniform way as above. We also give:

\begin{enumerate}
\def\labelenumi{(\Roman{enumi})}
\setcounter{enumi}{12}
\tightlist
\item
  The Cartan matrix, from which the Dynkin graph is derived as described in no. 2.
\end{enumerate}

